\section{Phương pháp hoạch định dựa trên đồ thị}

Các phương pháp hoạch định dựa trên đồ thị (Graph-based planning), tiêu biểu là thuật toán $A^*$ và Dijkstra, đóng vai trò nền tảng trong khoa học robot nhờ các đặc tính về tính tối ưu và tính đầy đủ theo độ phân giải (resolution completeness) trong không gian trạng thái được rời rạc hóa~\cite{LaValle2006}. Trong các kiến trúc phổ biến, hệ thống thường sử dụng bản đồ lưới chiếm chỗ (occupancy grid map) làm cơ sở dữ liệu để thuật toán $A^*$ thực hiện tìm kiếm quỹ đạo thông qua việc tối ưu hóa hàm chi phí kết hợp với hàm heuristic. Để hỗ trợ thuật toán hoạch định, hệ thống nhận thức thường tích hợp cảm biến LiDAR và khối đo lường quán tính (IMU). Một quy trình phổ biến là sử dụng thuật toán SLAM (như Cartographer) để xây dựng bản đồ và định vị đồng thời \cite{TandonEtAl2024}. Đối với các hệ thống hoạt động trong không gian 3D nhưng thực hiện hoạch định trong không gian trạng thái $SE(2)$, dữ liệu đám mây điểm (3D point cloud) từ cảm biến LiDAR thường được chiếu xuống mặt phẳng ngang nhằm đơn giản hóa mô hình môi trường trong khi vẫn duy trì các thông tin vật cản thiết yếu.

Tuy nhiên, sự hạn chế về tính liên tục của các phương pháp rời rạc thường dẫn đến các quỹ đạo thiếu độ trơn (zigzag), đòi hỏi các bước hậu xử lý phức tạp để thỏa mãn các ràng buộc động lực học của robot. Để giảm thiểu chi phí tính toán cho các hệ thống có cấu hình động học cao chiều, một chiến lược phổ biến là phân rã bài toán thành hai giai đoạn: hoạch định toàn cục (Global Planning) trong không gian tác vụ để xác định quỹ đạo thô, và thực thi cục bộ (Local Execution) để giải quyết các ràng buộc vi phân và duy trì sự ổn định của hệ thống~\cite{TandonEtAl2024}. Mặc dù cải thiện hiệu suất thời gian thực, sự tách biệt này thường tạo ra khoảng cách lớn giữa hình học và động lực học, dẫn đến các nghiệm dưới mức tối ưu hoặc mất tính khả thi trong các môi trường hẹp.

\section{Phương pháp hoạch định dựa trên lấy mẫu}

Để vượt qua ``lời nguyền đa chiều'' (curse of dimensionality) mà các phương pháp đồ thị truyền thống gặp phải, các thuật toán dựa trên lấy mẫu (SBP) như Probabilistic Roadmap Method (PRM) và Rapidly-exploring Random Trees (RRT) đã trở thành công cụ mạnh mẽ trong việc xử lý không gian cấu hình ($C$-space) số chiều cao~\cite{LaValle1998}. Thay vì xây dựng tường minh các rào cản vật lý, SBP xấp xỉ tính kết nối của không gian tự do ($C_{free}$) thông qua việc lấy mẫu ngẫu nhiên các cấu hình và kết nối chúng thành cấu trúc đồ thị hoặc cây. Trong khi PRM là bộ lập hoạch đa truy vấn dựa trên giai đoạn học và truy vấn bản đồ lộ trình, thì RRT lại tập trung vào đơn truy vấn bằng cách mở rộng cây từ điểm bắt đầu hướng tới mục tiêu thông qua cơ chế ``độ chệch Voronoi'' (Voronoi bias), hỗ trợ khám phá nhanh chóng các khu vực chưa xác định~\cite{LaValle1998}.

Mặc dù sở hữu tính đầy đủ theo xác suất (probabilistic completeness), các phương pháp SBP nguyên bản thường tạo ra quỹ đạo dạng ``răng cưa'' và thiếu tính tối ưu về mặt hình học. Biến thể RRT* đã giới thiệu các thủ tục tái kết nối (rewire) để đạt được tính tối ưu tiệm cận (asymptotic optimality), song tốc độ hội tụ của nó thường chậm và đòi hỏi tài nguyên tính toán lớn cho các bước làm mịn~\cite{GammellEtAl2014}. Thách thức lớn nhất đối với SBP vẫn nằm ở việc tích hợp các ràng buộc vi phân liên tục (kinodynamic constraints), do việc giải bài toán biên (BVP) để kết nối các mẫu rời rạc trong các hệ thống phi holonom là cực kỳ phức tạp và nhạy cảm với các hàm metric khoảng cách~\cite{GammellEtAl2014}.

\section{Tối ưu hóa quỹ đạo trực tiếp (Direct Trajectory Optimization - DTO)}

Các phương pháp tối ưu hóa quỹ đạo trực tiếp đóng vai trò quyết định trong việc tạo ra các chuyển động mượt mà và khả thi về mặt vật lý thông qua việc chuyển đổi bài toán điều khiển tối ưu liên tục thành bài toán lập trình phi tuyến (Nonlinear Programming - NLP)~\cite{Betts2010}. Trong nhóm phương pháp này, Direct Multiple Shooting chia quỹ đạo thành nhiều phân đoạn cục bộ, tích phân động lực học trên từng đoạn và áp đặt các ràng buộc ghép nối để bảo đảm tính liên tục của trạng thái tại các điểm nút~\cite{BockPlitt1984}. 

Tư tưởng chia bài toán đường đi thành các đoạn nhỏ và cập nhật các điểm bắn cũng xuất hiện trong bài toán đường đi ngắn nhất hình học. Cụ thể, Phan Thành An và Nguyễn Thị Lệ đề xuất một phương pháp multiple shooting để tìm đường đi xấp xỉ ngắn nhất cho robot tự hành trong môi trường 2D chưa biết, trong đó chuỗi các bó đoạn thẳng được phân hoạch, các shooting points được cập nhật lặp lại, và điều kiện thẳng hàng được dùng như tiêu chuẩn dừng~\cite{AnLe2023MMS}. Tuy nhiên, khác với hướng tiếp cận hình học này, tối ưu hóa quỹ đạo trực tiếp trong đồ án nhúng trực tiếp phương trình động lực học $\dot{x}=f(x,u)$ vào bài toán tối ưu dưới dạng các ràng buộc defect và ràng buộc ghép nối. Nhờ đó, phương pháp phù hợp hơn với các bài toán hoạch định chuyển động cần bảo đảm tính khả thi động lực học, thay vì chỉ tối ưu hình dạng đường đi trong không gian cấu hình.

Trong các nghiên cứu hiện đại, các bộ lập hoạch dựa trên lấy mẫu (sampling-based planners, SBP) thường được kết hợp với tối ưu hóa quỹ đạo trực tiếp như một bước hậu xử lý để tinh chỉnh nghiệm ban đầu, qua đó cân bằng giữa khả năng khám phá không gian và chất lượng quỹ đạo về năng lượng hoặc thời gian. Ở hướng tích hợp chặt chẽ hơn, các bộ lập hoạch dựa trên Lập trình lồi số nguyên hỗn hợp (Mixed-Integer Convex Programming, MICP) đã được đề xuất~\cite{MarcucciEtAl2023GCS}, với mục tiêu kết hợp tính đầy đủ theo xác suất của SBP và khả năng xử lý chính xác các ràng buộc động học, động lực học của robot.

